\section*{第 \ref{ch:b1:intermolecular-forces-solvents} 章——分子间作用力与溶剂}

\begin{solution}{exo:b1:intermolecular-forces-solvents:1}
氩：仅有\omterm{def:b1:intermolecular-forces-solvents:vdw}{色散力}。氯化氢：\omterm{def:b1:intermolecular-forces-solvents:vdw}{葛生力}、\omterm{def:b1:intermolecular-forces-solvents:vdw}{德拜力}和\omterm{def:b1:intermolecular-forces-solvents:vdw}{色散力}，
以\omterm{def:b1:intermolecular-forces-solvents:vdw}{色散力}为主（见\cref{exo:b1:intermolecular-forces-solvents:8}）；氯原子体积过大，
难以形成强\omterm{def:b1:intermolecular-forces-solvents:hbond}{氢键}。水：以\omterm{def:b1:intermolecular-forces-solvents:hbond}{氢键}为主，
并包含范德瓦耳斯力的三个分项。甲烷：仅有\omterm{def:b1:intermolecular-forces-solvents:vdw}{色散力}。二碘：仅有\omterm{def:b1:intermolecular-forces-solvents:vdw}{色散力}，
但因分子大、\omterm{def:b1:periodicity:polarisability}{极化率}高，作用较强。
\end{solution}

\begin{solution}{exo:b1:intermolecular-forces-solvents:2}
同种分子之间：\ce{CH3OH}（\ce{O-H}）、\ce{CH3NH2}
（\ce{N-H}）和 \ce{HF}。\ce{CH3OCH3} 与 \ce{CH3F} 的 O 或 F 上有\omterm{def:b1:lewis-resonance-vsepr:lewis}{孤电子对}，
但其上没有与之相连的氢原子，因此只能接受\omterm{def:b1:intermolecular-forces-solvents:hbond}{氢键}，例如来自水的\omterm{def:b1:intermolecular-forces-solvents:hbond}{氢键}。
\ce{CH4} 既不给出也不接受\omterm{def:b1:intermolecular-forces-solvents:hbond}{氢键}。
\end{solution}

\begin{solution}{exo:b1:intermolecular-forces-solvents:3}
水和乙醇：极性、质子性。丙酮和乙腈：极性、
非质子性。二氯甲烷：弱极性、非质子性。己烷：非极性、非质子性。
\end{solution}

\begin{solution}{exo:b1:intermolecular-forces-solvents:4}
二碘是非\omterm{def:b1:lewis-resonance-vsepr:dipole}{极性分子}：在己烷中，\omterm{def:b1:intermolecular-forces-solvents:vdw}{色散力}只是换成了另一种\omterm{def:b1:intermolecular-forces-solvents:vdw}{色散力}，
而在水中则会破坏\omterm{def:b1:intermolecular-forces-solvents:hbond}{氢键}，却无法加以替代。
氯化钠是离子化合物：只有\omterm{def:b1:intermolecular-forces-solvents:permittivity}{相对介电常数}高的溶剂
才能使其离子分开，只有\omterm{def:b1:intermolecular-forces-solvents:solvent-classes}{极性溶剂}才能使其\omterm{def:b1:intermolecular-forces-solvents:dissolution}{溶剂化}；
己烷（$\varepsilon_r = 1.89$）两者都做不到。
\end{solution}

\begin{solution}{exo:b1:intermolecular-forces-solvents:5}
二甲醚不含 \ce{O-H}，不能给出\omterm{def:b1:intermolecular-forces-solvents:hbond}{氢键}，因此沸点最低。
甲醇和乙醇都能形成\omterm{def:b1:intermolecular-forces-solvents:hbond}{氢键}；乙醇
多一个 \ce{CH2}，\omterm{def:b1:intermolecular-forces-solvents:vdw}{色散力}更强。乙酸形成\omterm{def:b1:intermolecular-forces-solvents:hbond}{氢键}二聚体，
相对分子质量最大，沸点最高。
\end{solution}

\begin{solution}{exo:b1:intermolecular-forces-solvents:6}
真空中 $F = e^2/(4\pi\varepsilon_0 r^2) = (1.602 \times
10^{-19})^2/(4\pi \times 8.854 \times 10^{-12} \times (5.00 \times
10^{-10})^2) = \qty{9.23e-10}{N}$。在水中 $F/78.4 = \qty{1.18e-11}{N}$；
在己烷中 $F/1.89 = \qty{4.88e-10}{N}$。水对吸引力的削弱，是己烷的 41 倍
：水能使离子保持分开，己烷则不能。
\end{solution}

\begin{solution}{exo:b1:intermolecular-forces-solvents:7}
丙酮与水混合时，会破坏部分水–水\omterm{def:b1:intermolecular-forces-solvents:hbond}{氢键}，但会与水形成新的\omterm{def:b1:intermolecular-forces-solvents:hbond}{氢键}，
即水与 \ce{C=O} 中的氧之间的\omterm{def:b1:intermolecular-forces-solvents:hbond}{氢键}，另有偶极–偶极作用，
总体上有利。己烷与水之间只能形成较弱的\omterm{def:b1:intermolecular-forces-solvents:vdw}{色散力}，
以换取其被破坏的\omterm{def:b1:intermolecular-forces-solvents:hbond}{氢键}，
因此两者保持分层。
\end{solution}

\begin{solution}{exo:b1:intermolecular-forces-solvents:8}
\ce{HBr} 的电子更多，电子云更大、更易极化，其
\omterm{def:b1:intermolecular-forces-solvents:vdw}{色散力}比 \ce{HCl} 的增量超过\omterm{def:b1:intermolecular-forces-solvents:vdw}{葛生力}的减少量，
因此以\omterm{def:b1:intermolecular-forces-solvents:vdw}{色散力}为主。
\end{solution}

\begin{solution}{exo:b1:intermolecular-forces-solvents:9}
\chemfig{CH_3-C(=[:60]O)-[:-60]O-[:0]H} 与其同类相对，两个
$\ce{O-H}\cdots\ce{O=C}$ 键闭合成一个八元环。在苯中，
非极性、非质子，无物质与这些\omterm{def:b1:intermolecular-forces-solvents:hbond}{氢键}竞争，乙酸
发生二聚化，表现为约 \qty{120}{g/mol} 的粒子。在水中，
每个乙酸分子改与水形成\omterm{def:b1:intermolecular-forces-solvents:hbond}{氢键}，二聚体因而解离。
\end{solution}

\begin{solution}{exo:b1:intermolecular-forces-solvents:10}
氯化钠由离子构成：\omterm{def:b1:intermolecular-forces-solvents:permittivity}{相对介电常数}高的水能使其分开（解离），
并使离子水合（氧原子朝向 \ce{Na+}，氢原子
朝向 \ce{Cl-}）。氯化氢为\omterm{def:b1:lewis-resonance-vsepr:dipole}{极性分子}：极性的水
使 \ce{H-Cl} 键转变为离子对 \ce{H3O+ Cl-}（电离），
再使离子对分开（解离），并使离子水合。己烷为非极性、
\omterm{def:b1:intermolecular-forces-solvents:permittivity}{相对介电常数}低，既不电离也不解离，
因此没有离子来传导电流。
\end{solution}

\begin{solution}{exo:b1:intermolecular-forces-solvents:11}
$(174.1 - 36.1)/5 = \qty{27.6}{\celsius}$，即每个 \ce{CH2} 的平均增量。
最后一个增量（壬烷到癸烷）为 \qty{23.4}{\celsius}，
因此十一烷的沸点应接近 $174 + 22 \approx \qty{196}{\celsius}$。
每个 \ce{CH2} 贡献的\omterm{def:b1:intermolecular-forces-solvents:vdw}{色散力}相同，
但在分子总体中所占比例越来越小，
而克服分子间作用所需的温度已经很高。
% ledger: bp:C9H20
\end{solution}

\begin{solution}{exo:b1:intermolecular-forces-solvents:12}
亲水部分为硫酸根头基 \ce{-OSO3^-} 及其抗衡离子 \ce{Na+}；
疏水部分为十二碳链。在表面，头基朝向水，尾链朝向空气；
在胶束中，尾链在内，头基在外。
尾链溶于污渍中的油脂，头基使其与水接触，
油脂便以胶束的形式脱离。
\end{solution}

\begin{solution}{pb:b1:intermolecular-forces-solvents:1}
\textbf{1.} \omterm{def:b1:intermolecular-forces-solvents:vdw}{色散力}是非极性原子之间唯一的作用力。
\textbf{2.} 从氦到氙，电子数和原子半径增大，
因此\omterm{def:b1:periodicity:polarisability}{极化率}和\omterm{def:b1:intermolecular-forces-solvents:vdw}{色散力}增强。
\textbf{3.} 氦有两个电子，被原子核紧紧束缚，
是\omterm{def:b1:periodicity:polarisability}{极化率}最小的原子，\omterm{def:b1:intermolecular-forces-solvents:vdw}{色散力}也最弱。
\textbf{4.} $68.8 - 36.1 = \qty{32.7}{\celsius}$：每增加一个 \ce{CH2}
增加电子数和接触面，因而\omterm{def:b1:intermolecular-forces-solvents:vdw}{色散力}增强。
\textbf{5.} \omterm{def:b1:intermolecular-forces-solvents:vdw}{色散力}随\omterm{def:b1:periodicity:polarisability}{极化率}增大而增强：$\ce{CH4} < \ce{SiH4} <
\ce{GeH4}$。
\textbf{6.} 碳的\omterm{def:b1:periodicity:electronegativity}{电负性}不够大，不能使 \ce{C-H} 键强烈极化，
且甲烷没有可接受\omterm{def:b1:intermolecular-forces-solvents:hbond}{氢键}的\omterm{def:b1:lewis-resonance-vsepr:lewis}{孤电子对}。
\textbf{7.} 水：极性、质子性，\omterm{def:b1:intermolecular-forces-solvents:permittivity}{相对介电常数}最高；它
既能使两种离子分开，又能使其\omterm{def:b1:intermolecular-forces-solvents:dissolution}{溶剂化}。
\textbf{8.} $78.4/1.89 = 41$：水中的吸引力比己烷中弱 41 倍。
\textbf{9.} 乙腈（或丙酮）极性足以溶解该盐，
且为\omterm{def:b1:intermolecular-forces-solvents:solvent-classes}{非质子性溶剂}，使阴离子保持游离状态。
\textbf{10.} 己烷与水不互溶，与二碘一样为非极性。碘
进入己烷层，即上层（己烷密度小于水）。
\textbf{11.} 乙醇与水互溶，不会形成第二层。
\textbf{12.} 其氧原子可接受来自水的\omterm{def:b1:intermolecular-forces-solvents:hbond}{氢键}，但自身不能给出\omterm{def:b1:intermolecular-forces-solvents:hbond}{氢键}，
且两个乙基是\omterm{def:b1:intermolecular-forces-solvents:amphiphile}{疏水的}，
因此仅微溶。
\textbf{13.} $-66.4 - (-85.0) = \qty{18.6}{\celsius}$，每\omterm{def:b1:periodicity:table}{周期}；
外推得 \ce{HF}：$-85.0 - 18.6 = \qty{-103.6}{\celsius}$。
\textbf{14.} \ce{HF} 的沸点为 \qty{19.6}{\celsius}，
比外推值高 \qty{123}{\celsius}。
\textbf{15.} $-62.5 - (-87.8) = \qty{25.3}{\celsius}$；外推得
\ce{NH3}：\qty{-113.1}{\celsius}；实际值为 $-33.4$，
高出 \qty{80}{\celsius}。
\textbf{16.} $-88.1 - (-112.0) = \qty{23.9}{\celsius}$；外推得
\ce{CH4}：\qty{-135.9}{\celsius}；实际值为 $-161.5$，低于外推值。甲烷
不形成\omterm{def:b1:intermolecular-forces-solvents:hbond}{氢键}，仅因其分子小、\omterm{def:b1:periodicity:polarisability}{极化率}低。
\textbf{17.} 水（约 \qty{179}{\celsius}）$>$ \ce{HF}（123）$>$
\ce{NH3}（80）。
\textbf{18.} \ce{NH3} 可给出三个\omterm{def:b1:intermolecular-forces-solvents:hbond}{氢键}，但只能接受一个（仅一个\omterm{def:b1:lewis-resonance-vsepr:lewis}{孤电子对}）；
\ce{HF} 可接受三个\omterm{def:b1:intermolecular-forces-solvents:hbond}{氢键}，但只能给出一个。水既给出两个，又接受两个，
每个分子可参与四个\omterm{def:b1:intermolecular-forces-solvents:hbond}{氢键}，
形成完整的三维网络。
\textbf{19.} \ce{H2S} 与 \ce{H2Se} 不形成明显的\omterm{def:b1:intermolecular-forces-solvents:hbond}{氢键}：
其沸点遵循\omterm{def:b1:intermolecular-forces-solvents:vdw}{色散力}决定的趋势，
若无\omterm{def:b1:intermolecular-forces-solvents:hbond}{氢键}，水也将遵循此趋势。水本身才是此处的反常值。
\textbf{20.} $-41.3 - (-60.3) = \qty{19.0}{\celsius}$，每\omterm{def:b1:periodicity:table}{周期}。
\textbf{21.} $T(p) = -60.3 + 19.0\,(p - 3)$，单位为 \unit{\celsius}。
\textbf{22.} $T(2) = -60.3 - 19.0 = \qty{-79.3}{\celsius}$。
\textbf{23.} $100.0 - (-79.3) = \qty{179}{\celsius}$。
\textbf{24.} 若没有\omterm{def:b1:intermolecular-forces-solvents:hbond}{氢键}，水的沸点将约为
\textbf{\qty{-79}{\celsius}}，在地球上将呈气态。
\end{solution}
